\documentclass{article}
\usepackage{graphicx} % Required for inserting images
\usepackage{geometry}
\usepackage[style=authoryear, backend=biber]{biblatex}
\addbibresource{bibliography.bib}
\usepackage{hyperref}
\usepackage{booktabs}
\usepackage{tabularx}
\usepackage{url}
\title{Reducing Online Extremism with a [Pre-]De-radicalizing Language Agent\\\Large{Deliverable 1 Report}}
\author{Andrew Shaw, Advait Bhat, Aahan Patel,\\Galen Weld, Amy X. Zhang, Tim Althoff}
\date{March 2026}

\begin{document}

\maketitle

\section{Introduction}\label{sec:intro}

\section{Background \& Related Work}\label{sec:related}

\subsection{Radicalization and Polarization}\label{sec:radicalization_polarization}
Whereas polarization is typically defined as ``the process through which people extremize in their views through communication with likeminded others about their grievances,'' radicalization can be understood as ``a form of polarization that involves internalizing specific norms for [...] violent collective action'' \parencite{Smith2025PrimerOP}. The contributing factors to radicalization are complex and debated, but are generally understood to include the experience of a grievance, social bonds with others who hold radical beliefs, the extremist ideology itself, and a supporting environment for radicalization to occur \parencite{Leap2021RadicalizationAD}. Moreover, radicalizaton is most likely to occur when groups perceive that conventional political strategies are unable to create change, meaning that restrictive policies which close off paths for legitimate contestation can paradoxically end up fueling greater radicalization \parencite{Smith2025PrimerOP}.

One challenge in responding to radicalization as defined above is that the line between ``mere'' polarization and radicalization may not always be clear, since the distinction between the two hinges on the degree to which polarized behavior also reflects violent belief. Compounding this challenge is the fact that polarization may play an essential role in forming political identities and fostering political engagement, which are themselves important features of healthy pluralistic and democratic societies \parencite{Smith2024PolarizationIT}. Indeed, given our limited ability to process the sheer volume of information available in today's online ecosystems, polarization can be considered the result of an entirely rational decision to filter information through one's trusted epistemic community \parencite{Almagro2022PoliticalPR}. A useful framework for conceptualizing both the dangers and benefits of polarization is  \textit{agonistic pluralism}, a democratic theory which raises a similar distinction between ``antagonistic'' and ``agonistic'' political conflict. According to agonistic theorists like \textcite{Mouffe2000DemocraticParadox}, participants of \textit{antagonistic} conflict view opponents as an ``enemy to be destroyed,'' while participants of \textit{agonistic} conflict view opponents as an adversary ``whose ideas we combat but whose right to defend those ideas we do not put into question.'' Importantly, this does not entail that antagonistic views are always unacceptable to hold. Rather, \textcite{Westphal2024AffectivePA} suggests that antagonistic conflict can be tolerated in a pluralistic society so long as it does not \textit{remain} antagonistic; the task of democratic politics is thus to provide mechanisms for defusing antagonistic relationships into agonistic ones.

\subsection{AI and Political Discourse}\label{sec:ai_political_discourse}
Recent scholarship has increasingly taken note of the potential for AI technologies like LLMs to influence political discourse \parencite{Fisher2025BiasedLLMs}, but remains split on how these technologies may affect phenomena like polarization and radicalization. Some scholars like \textcite{Hendricks2026LLMsCF} have raised concerns that LLMs can fuel extremist attitudes in conversation participants by using universal moral framings. In contrast, in a different experiment, \textcite{Sevi2025AIWB} find that conversations with LLMs only weaken political attitudes, but do not strengthen them. Other scholars have found more proactive LLM-based interventions to be effective at reducing conspiracy beliefs \parencite{Costello2024DurablyRC}, helping participants find common ground \parencite{Tessler2024HabermasMachine}, promoting feelings of being understood in online conversations \parencite{Argyle2023LeveragingAI}, and otherwise reducing polarization.

One approach towards reconciling these divergent findings is to look to individual-level differences. Earlier work from \textcite{Munson2010PresentingDP} finds that news readers broadly exhibit two types of behavior: diversity-seeking and challenge-averse. Whereas diversity-seeking readers report higher levels of satisfaction when presented with diverse political opinions, while challenge-averse readers report the opposite. Similarly, \textcite{Rathje2025SycophanticAI} find that users with higher self-reported open-mindedness are more receptive to disagreeable chatbots, while users with lower self-reported open-mindedness prefer validating chatbots. They therefore suggest that fostering open-mindedness and presenting counter-attitudinal facts in a validating way could both be strategies for increasing engagement with diverse perspectives.

\textcite{Lazar2026UsingLL} carry out a comprehensive evaluation of depolarizing applications of LLMs, along with broader proposals of using LLMs to enhance democracy. They find that these projects generally leverage one of four strategies: using LLMs to (1) \textit{summarize} public submissions to government consultations, (2) generate widely acceptable statements that \textit{aggregate} diverse political opinions, (3) \textit{represent} individual preferences over unseen policy options, and (4) \textit{facilitate} democratic deliberation. However, they note that to date, many such projects have been carried out under experimental conditions that closely resemble the ``ideal speech situation'' proposed by \textcite{Habermas1985TheoryOC} in which participants engage in rational, unmanipulated, and dispassionate debate with others---with some proposals explicitly citing Habermas' work as inspiration for their design \parencite{Tessler2024HabermasMachine}. It therefore remains unclear how well these proposals will generalize to the everyday settings in which political discourse is most often carried out, structured by inequalities, deep ideological disagreement, economic and social pressures, and adversarial relationships. Indeed, it is on this same point that \textcite{Mouffe2000DemocraticParadox} objects to the deliberative democratic notion of a rational consensus as an illusory abstraction. By pointing instead to the necessity of learning to constructively negotiate across conflict and difference, agonistic pluralism may here offer a helpful framework for designing for such non-ideal speech situations.



\subsection{Community Notes}\label{sec:comm_notes}
Community notes, originally called `Birdwatch,' was developed by Twitter (now X) and launched in January 2021~\cite{Wojcik2022BirdwatchCW}.

\section{Proposed Experimental Design}\label{sec:design}
To realistically evaluate the efficacy of our de-radicalizing interventions, we will deploy our tool in a randomized controlled trial on Twitter (now X), one of the largest social media platforms in the world. Our participants will be recruited to serve in a `bystander' role, invoking the tool when they encounter Tweets made by potentially radicalized users.

When a \textbf{bystander participant} `@-mentions' the tool to invoke it in reply to a \textbf{target tweet}, our system randomly assigns a treatment condition (described in Section~\ref{sec:conditions}).

\subsection{Treatment Conditions}\label{sec:conditions}

\subsection{Outcome Measurements}\label{sec:outcomes}



\section{Pilot Results \& System Functionality}\label{sec:demo}

We have implemented a working prototype of our intervention system as a retrieval-augmented generation (RAG) agent grounded in Twitter's Community Notes dataset. This section describes the current system design (its a v1 prototype and will be updated in the future), explains how the two treatment conditions are operationalised in this prototype, and presents pilot outputs demonstrating functionality across both conditions.

\subsection{System Design}\label{sec:design}

We have implemeneted a simplified version 1 of our final pipeline. This prototype focuses on implementing the retrieval and generation functionality of our system, while deffering the implementation of user polarity scoring (i.e. placing the user we are replying to on a left-right political spectrum using community notes data) for the nest iteration.

The current system works as follows: Given a target tweet, our system first uses an LLM to decompose the claim into several focused search queries, which are then embedded and matched against a FAISS vector index~\cite{Johnson2019FAISS} built from approximately 2.4 million publicly available Community Notes. For each retrieved seed note, all other notes written about the same tweet are pulled in, capturing the full deliberative context around that piece of content. The resulting set is deduplicated and each note is scored on a continuous misleadingness axis using labels in the dataset (Community Notes classifications and rater flags). Notes are then sampled proportionally across misleadingness buckets so that the evidence presented to the response model reflects the actual distribution of community assessments.

The sampled notes are passed to an LLM as private background knowledge. The model is instructed to write a direct reply to the original post informed by the evidence but without referencing notes, while using sources from the notes as references. The entire retrieval and scoring pipeline is shared across treatment conditions; the only variable is the system prompt (preamble) in the final response prompt.

\subsection{Treatment Conditions}\label{sec:condition_impl}

Our two conditions share the same retrieved evidence, grounding rules, and output schema, ensuring that observed differences in participant reactions can be attributed to the response's framing rather than to differences in the underlying evidence.

The \textbf{bridging} condition instructs the model to acknowledge what is understandable about the poster's concern before presenting corrective information, and to use inclusive, calm language that frames the reply as 'adding context' rather than 'correcting the person'.

The \textbf{agonistic} condition instructs the model to engage with the claim as a serious assertion that deserves a rebuttal by stating disagreement clearly and presenting counter-evidence directly. This results in the model being direct and not softening its rebuttal.

\subsection{Pilot Examples}\label{sec:pilot_examples}

To demonstrate system functionality, we ran the pipeline against two claims representative of common polarization topics, generating responses under both treatment conditions. Table~\ref{tab:pilot_examples} presents two illustrative examples.

\begin{table}[t]
\centering
\footnotesize
\caption{Pilot outputs for two polarizing claims under both treatment conditions.}
\label{tab:pilot_examples}
\begin{tabularx}{\columnwidth}{>{\raggedright\arraybackslash}X >{\raggedright\arraybackslash}X}
\toprule
\textbf{Bridging} & \textbf{Agonistic} \\
\midrule
\multicolumn{2}{c}{\textit{``Vaccines cause autism.''}} \\
\midrule
I get why that worry spreads---hearing a clear cause for something as painful as autism feels like an explanation. Scientists who study this have not found vaccines to be the cause, and the early paper that started the fear was later found to be fraudulent and discredited. Concerns about a mercury preservative led to its removal from most childhood vaccines years ago, yet autism diagnoses did not change in a way that supports a link. The best-supported conclusion is that vaccines are not the cause of autism. \newline\textbf{Sources:} \href{https://www.cdc.gov/vaccinesafety/concerns/autism.html}{CDC}, \href{https://www.nationalacademies.org/based-on-science/vaccines-do-not-cause-autism}{Nat.\ Academies}, \href{https://www.factcheck.org/2023/07/scicheck-false-claim-about-cause-of-autism-highlighted-on-pennsylvania-senate-panel/}{FactCheck.org}
& I disagree---vaccines do not cause autism. The original claim came from a fraudulent study that was retracted, and many well-conducted studies since then have found no causal link. Thimerosal, the preservative often blamed, was removed from routine childhood vaccines years ago yet autism diagnoses did not track with that change. Vaccination protects children from serious diseases and the autism claim is not supported by reliable research. \newline\textbf{Sources:} \href{https://ncbi.nlm.nih.gov/pmc/articles/PMC2831678/}{NCBI}, \href{https://acpjournals.org/doi/10.7326/M18-2101}{Annals of Internal Med.}, \href{https://cdc.gov/mmwr/volumes/65/ss/ss6503a1.htm}{CDC MMWR} \\
\midrule
\multicolumn{2}{c}{\textit{``The 2020 american presidential election was stolen through mass voter fraud.''}} \\
\midrule
I get why you'd feel strongly about that---trust in elections matters. The legal challenges brought after 2020 were repeatedly dismissed by judges, and senior election officials at the time said the process was secure. While no system is perfect and isolated fraud can happen in any election, there isn't credible support that the 2020 result was overturned by mass, coordinated voter fraud. \newline\textbf{Sources:} \href{https://www.reuters.com/article/world/fact-check-courts-have-dismissed-multiple-lawsuits-of-alleged-electoral-fraud-p-idUSKBN2AF1FQ/}{Reuters}, \href{https://campaignlegal.org/results-lawsuits-regarding-2020-elections}{Campaign Legal Center}, \href{https://www.brookings.edu/articles/widespread-election-fraud-claims-by-republicans-dont-match-the-evidence/}{Brookings}
& That's not true. Courts across multiple states reviewed the fraud claims and rejected them, often for lack of proof, and the official, certified results in the key battleground states show Joe Biden won those contests. Independent reviews found no widespread or systemic voter fraud that would have overturned the outcome; isolated irregularities do not equal a stolen election. \newline\textbf{Sources:} \href{https://www.reuters.com/article/world/fact-check-courts-have-dismissed-multiple-lawsuits-of-alleged-electoral-fraud-p-idUSKBN2AF1FQ/}{Reuters}, \href{https://campaignlegal.org/results-lawsuits-regarding-2020-elections}{Campaign Legal Center}, \href{https://apnews.com/article/voting-rights-politics-united-states-presidential-elections-denver-40238a64632af33f5ded28391ed83a8a}{AP News} \\
\bottomrule
\end{tabularx}
\end{table}

\printbibliography

\end{document}
